% Part: normal-modal-logic
% Chapter: axioms-systems
% Section: soundness

\documentclass[../../../include/open-logic-section]{subfiles}

\begin{document}

\olfileid[id]{nml}{prf}{snd}

\olsection{Kesahihan}

Suatu sistem !!{derivation} disebut sahih jika semua yang dapat !!{derive}
bersifat valid. Ketika membahas sistem modal, yakni !!{derivation}s yang
selain menggunakan~\Ax{K} juga dapat menggunakan instansiasi beberapa
!!{formula}s $!A_1$, \dots,~$!A_n$, kita menginginkan agar setiap formula
yang !!{derivable} benar dalam setiap model tempat $!A_1$, \dots, $!A_n$
benar.

\begin{thm}[Teorema Kesahihan]\ollabel{thm:soundness}
  Jika setiap instansiasi dari $!A_1$, \dots, $!A_n$ valid, berturut-turut,
  dalam kelas model $\mClass{C}_1$, \dots, $\mClass{C}_n$, maka
  $\Log{K}!A_1\dots !A_n\Proves !B$ mengimplikasikan bahwa $!B$ valid dalam
  kelas model $\mClass{C}_1\cap\dots\cap\mClass{C}_n$.
\end{thm}

\begin{proof}
  Melalui induksi pada panjang pembuktian. Agar ringkas, tetapkan
  $\mClass{C}=\mClass{C}_1\cap\dots\cap\mClass{C}_n$.
  \begin{enumerate}
  \item Dasar induksi: Jika $!B$ mempunyai pembuktian dengan panjang~$1$,
    maka formula itu merupakan instansiasi tautologis, instansiasi
    \Ax{K},\iftag{prvDiamond}{ atau~\Dual{},}{} atau instansiasi salah satu
    dari $!A_1$, \dots,~$!A_n$. Dalam kasus pertama, $!B$ valid dalam
    $\mClass{C}$, sebab instansiasi tautologis valid dalam kelas model
    \emph{mana pun}, menurut \olref[syn][tau]{prop:valid-taut}. Hal yang sama
    berlaku dalam kasus kedua menurut
    \olref[syn][sch]{prop:Kvalid}\iftag{prvDiamond}{ dan
      \olref[syn][sch]{prop:Dual-valid}}{}. Terakhir, dalam kasus ketiga,
    karena $!B$ valid dalam $\mClass{C}_i$ dan
    $\mClass{C}\subseteq\mClass{C}_i$, $!B$ juga valid dalam~$\mClass{C}$
    menurut \olref[syn][val]{prop:subset-class}.
  \item Langkah induksi: Andaikan $!B$ mempunyai pembuktian dengan panjang
    $k>1$. Jika $!B$ merupakan instansiasi tautologis, instansiasi
    \Ax{K},\iftag{prvDiamond}{ instansiasi~\Dual{},}{} atau instansiasi salah
    satu dari $!A_1$, \dots, $!A_n$, kita melanjutkan seperti pada langkah
    sebelumnya. Jadi, andaikan $!B$ diperoleh melalui~\MP{} dari !!{formula}s
    sebelumnya $!C\lif !B$ dan~$!C$. Maka $!C\lif !B$ dan~$!C$ mempunyai
    pembuktian dengan panjang~$<k$, dan menurut hipotesis induksi keduanya
    valid dalam~$\mClass{C}$. Menurut \olref[syn][sch]{prop:soundMP}, $!B$
    juga valid dalam~$\mClass{C}$. Terakhir, andaikan $!B$ diperoleh melalui
    \Nec{} dari~$!C$ (sehingga $!B=\Box!C$). Menurut hipotesis induksi, $!C$
    valid dalam~$\mClass{C}$; menurut \olref[syn][val]{prop:Nec-rule},
    demikian pula~$!B$.
  \end{enumerate}
\end{proof}

\end{document}
